% Research notes for the parent report.  This is an input fragment.
% Standard amsmath, amssymb, amsthm theorem/proof environments are assumed.

\section{A Barnes-integral companion}
\label{bim:section}

This section records evaluated consequences of the classical first Barnes
lemma.  The underlying integral is not new: it is NIST DLMF
\href{https://dlmf.nist.gov/5.13.E3}{Equation 5.13.3} \cite{DLMF}.  The present
contribution is the selection and complete differentiation of the
centered fourth moment and the pure trigamma correlation.  No historical priority claim is made for these specializations.  These ordinary integrals are
distinct from the higher-dimensional Gamma density used in the recent
\emph{Gamma Normalizations and Colored Moments} report.

For complex parameters of positive real part, put
\begin{align}
 W(a,b,c,d;x)&=\Gamma(a+ix)\Gamma(b+ix)
                  \Gamma(c-ix)\Gamma(d-ix),\\
 F(a,b,c,d)&=\frac{\Gamma(a+c)\Gamma(a+d)
                       \Gamma(b+c)\Gamma(b+d)}
                      {\Gamma(a+b+c+d)}.
\end{align}
Barnes' first lemma states
\begin{equation}\label{bim:barnes}
 \frac1{2\pi}\int_{-\infty}^{\infty}W(a,b,c,d;x)\,dx
                              =F(a,b,c,d).
\end{equation}
Every finite parameter derivative can be passed under this integral.
Indeed, on a compact parameter set inside the four right half-planes,
Stirling's formula gives exponential decay $e^{-2\pi|x|}$ times a
fixed polynomial in $1+|x|$; differentiating adds only a fixed power of
$1+\log(2+|x|)$.  This also proves that all integrals below are ordinary
absolutely convergent integrals.

\subsection{Centered mixed moments}

For $\sigma>0$ define
\begin{equation}\label{bim:density}
 p_\sigma(x)=\frac{\Gamma(4\sigma)}
                   {2\pi\Gamma(2\sigma)^4}
                        |\Gamma(\sigma+ix)|^4,
 \qquad \int_{\mathbb R}p_\sigma(x)\,dx=1.
\end{equation}
Write $\mathbb E_\sigma$ for integration against this density, and put
\begin{equation}\label{bim:center}
 m_\sigma=2\psi(2\sigma)-\psi(4\sigma),\qquad
 U_\sigma(x)=\psi(\sigma+ix)-m_\sigma.
\end{equation}
The following formulas involve only ordinary polygamma values:
\begin{align}
 \mathbb E_\sigma U_\sigma&=0,\label{bim:mean}\\
 \mathbb E_\sigma U_\sigma^2&=-\psi'(4\sigma),
                                               \label{bim:square}\\
 \mathbb E_\sigma|U_\sigma|^2
          &=\psi'(2\sigma)-\psi'(4\sigma),      \label{bim:abs-square}\\
 \mathbb E_\sigma\bigl(U_\sigma|U_\sigma|^2\bigr)
          &=-\psi''(4\sigma),                  \label{bim:cubic}\\
 \mathbb E_\sigma|U_\sigma|^4
          &=\psi'(4\sigma)^2
             +2\bigl(\psi'(2\sigma)-\psi'(4\sigma)\bigr)^2
             -\psi'''(4\sigma).                \label{bim:fourth}
\end{align}
The left side of \eqref{bim:cubic} is real by the change $x\mapsto-x$.

\begin{proof}
Set $L=\log F$.  At $a=b=c=d=\sigma$, all first logarithmic
derivatives equal $m_\sigma$.  Subtract this linear term by replacing
$F/F(\sigma,\sigma,\sigma,\sigma)$ with
\[
 \widetilde F(a,b,c,d)
 =e^{-m_\sigma(a+b+c+d-4\sigma)}
      \frac{F(a,b,c,d)}{F(\sigma,\sigma,\sigma,\sigma)}.
\]
Its single-slot integral insertions are $U_\sigma$ in the slots $a,b$
and $\overline U_\sigma$ in $c,d$.  The first derivatives of
$\log\widetilde F$ vanish.  For distinct slots the remaining derivatives
needed here are
\begin{align*}
 L_{ab}=L_{cd}&=q=-\psi'(4\sigma),\\
 L_{ac}=L_{ad}=L_{bc}=L_{bd}
      &=p=\psi'(2\sigma)-\psi'(4\sigma),\\
 L_{abc}=L_{abd}=L_{acd}=L_{bcd}&=-\psi''(4\sigma),\\
 L_{abcd}&=-\psi'''(4\sigma).
\end{align*}
The derivatives of $\widetilde F=e^{\log\widetilde F}$ of orders two
and three give the first four assertions.  Its $abcd$ derivative is
$L_{ab}L_{cd}+L_{ac}L_{bd}+L_{ad}L_{bc}+L_{abcd}$, because every
partition containing a singleton has zero contribution.  This is
$q^2+2p^2-\psi'''(4\sigma)$ and proves \eqref{bim:fourth}.
\end{proof}

\subsection{Eliminating all digamma insertions from a trigamma correlation}

\begin{theorem}\label{bim:trigamma}
For $\sigma>0$ one has
\begin{align}
 \mathbb E_\sigma\psi'(\sigma+ix)&=2\psi'(2\sigma),
                                                   \label{bim:trigamma-mean}\\
 \mathbb E_\sigma\left|\psi'(\sigma+ix)\right|^2
                  &=4\psi'(2\sigma)^2+\psi'''(2\sigma).
                                                   \label{bim:trigamma-square}
\end{align}
More generally, if $\Re u,\Re v>0$ and $s=u+v$, then
\begin{equation}\label{bim:unequal-trigamma}
\begin{split}
 &\int_{\mathbb R}\Gamma(u+ix)^2\Gamma(v-ix)^2
                         \psi'(u+ix)\psi'(v-ix)\,dx\\
 &\qquad=2\pi\frac{\Gamma(s)^4}{\Gamma(2s)}
                           \bigl(4\psi'(s)^2+\psi'''(s)\bigr).
\end{split}
\end{equation}
\end{theorem}

\begin{proof}
Use the constant-coefficient operators
\[
 D_+=\partial_a^2-\partial_a\partial_b,
 \qquad D_-=\partial_c^2-\partial_c\partial_d.
\]
After setting $a=b=u$ the first operator inserts exactly
$\psi'(u+ix)$, since the two copies of $\psi(u+ix)^2$ cancel.  The
second operator inserts $\psi'(v-ix)$ after $c=d=v$.

For an explicit right-side calculation, center as above, now at
$a=b=u,c=d=v$; the common first logarithmic derivative is
$2\psi(s)-\psi(2s)$.  Put
\[
 \begin{gathered}
 A=\psi'(s),\quad B=\psi'(2s),\quad
 A_3=\psi'''(s),\quad B_3=\psi'''(2s),\\
 d=2A-B,\quad q=-B,\quad p=A-B.
 \end{gathered}
\]
Centering does not affect $D_+D_-\widetilde F$ at coalescence: the extra
terms contain $\partial_a-\partial_b$ or
$\partial_c-\partial_d$, which vanish by the respective symmetries.
The required derivatives are
\begin{align*}
 \widetilde F_{aacc}&=d^2+2p^2+A_3-B_3,\\
 \widetilde F_{aacd}=\widetilde F_{abcc}&=dq+2p^2-B_3,\\
 \widetilde F_{abcd}&=q^2+2p^2-B_3.
\end{align*}
Consequently
\[
 D_+D_-\widetilde F=(d-q)^2+A_3=4A^2+A_3.
\]
Multiplying by $2\pi F(u,u,v,v)$ proves
\eqref{bim:unequal-trigamma}; setting $u=v=\sigma$ proves
\eqref{bim:trigamma-square}.  The second-order calculation
$\widetilde F_{aa}-\widetilde F_{ab}=d-q=2A$ gives
\eqref{bim:trigamma-mean}.
\end{proof}

The unequal-shift formula can equally be written using
$\psi'(z)=\zeta(2,z)$ and $\psi'''(z)=6\zeta(4,z)$.  Thus it is a
Gamma-weighted Hurwitz-zeta convolution with two independently chosen
argument shifts.  The only shift invariant on the right is their sum.

\subsection{Explicit half- and quarter-argument evaluations}

At $\sigma=1/2$, the Gamma reflection formula makes
$p_{1/2}(x)=(\pi/2)\operatorname{sech}^2(\pi x)$ and
$U_{1/2}(x)=\psi(1/2+ix)+\gamma+1$.  Equations
\eqref{bim:abs-square}, \eqref{bim:fourth}, and
\eqref{bim:trigamma-square} become
\begin{align}
 \int_{\mathbb R}\frac{|\psi(1/2+ix)+\gamma+1|^2}
                            {\cosh^2(\pi x)}\,dx
                   &=\frac2\pi,\label{bim:half-square}\\
 \int_{\mathbb R}\frac{|\psi(1/2+ix)+\gamma+1|^4}
                            {\cosh^2(\pi x)}\,dx
                   &=\frac2\pi\left(9-\frac{\pi^2}{3}
                                          -\frac{7\pi^4}{180}\right),
                                               \label{bim:half-fourth}\\
 \int_{\mathbb R}\frac{|\psi'(1/2+ix)|^2}
                            {\cosh^2(\pi x)}\,dx
                   &=\frac{16\pi^3}{45}.\label{bim:half-trigamma}
\end{align}
Here the cancellation of Euler's constant is exact after centering.

At $\sigma=1/4$, put $V(x)=\psi(1/4+ix)+\gamma+4\log2$.
The normalization is $p_{1/4}(x)=|\Gamma(1/4+ix)|^4/(2\pi^3)$, and
one obtains the equally short evaluations
\begin{align}
 \int_{\mathbb R}|\Gamma(1/4+ix)|^4|V(x)|^2\,dx&=\frac{2\pi^5}{3},
                                                   \label{bim:quarter-square}\\
 \int_{\mathbb R}|\Gamma(1/4+ix)|^4|V(x)|^4\,dx&=\frac{11\pi^7}{30},
                                                   \label{bim:quarter-fourth}\\
 \int_{\mathbb R}|\Gamma(1/4+ix)|^4|\psi'(1/4+ix)|^2\,dx&=4\pi^7.
                                                   \label{bim:quarter-trigamma}
\end{align}
The special-value reductions use DLMF
\href{https://dlmf.nist.gov/5.15.E2}{5.15.2} and
\href{https://dlmf.nist.gov/5.15.E3}{5.15.3}.

